\documentclass[11pt]{article}
\usepackage[margin=1in]{geometry}
\usepackage{amsmath,amssymb}
\usepackage{booktabs}
\usepackage{hyperref}

\title{FPChecker Tutorial Example 6: Mathematical Formulation and Mixed-Precision Algorithmic Changes}
\author{}
\date{\today}

\begin{document}
\maketitle

\section{Purpose}
This document describes the mathematics implemented by the three programs in Example 6:
\begin{itemize}
  \item \texttt{fp32.cpp}: baseline FP32 algorithm,
  \item \texttt{mixed.cpp}: report-driven mixed-precision refinement,
  \item \texttt{fp64.cpp}: FP64 reference.
\end{itemize}

The goal is to show how high-relative-error lines in FP32 motivate both precision promotion and algorithmic changes.

\section{Input and Generated Data}
Given inputs
\[
N,\quad A,\quad a,\ b,\ c,
\]
the synthetic data sequence is
\[
x_i = T + A\sin(0.01 i) + 0.25(i \bmod 9), \quad i=0,\dots,N-1,
\]
with trend constant
\[
T = 10^8.
\]

In the code, \(A\) is \texttt{amp}. In \texttt{fp32.cpp} and \texttt{mixed.cpp}, the stored array uses \texttt{float}; in \texttt{fp64.cpp}, it uses \texttt{double}.

\section{Quantity 1: Small Root of a Quadratic}
The quadratic equation is
\[
a r^2 + b r + c = 0.
\]

\subsection{FP32 baseline (unstable form)}
The FP32 program computes
\[
\Delta = b^2 - 4ac,
\]
\[
r_{\text{small}}^{(\mathrm{fp32})} = \frac{-b + \sqrt{\Delta}}{2a}.
\]

When \(b \gg c\) and \(a\approx 1\), we get \(\sqrt{\Delta} \approx b\), so
\[
-b + \sqrt{\Delta}
\]
subtracts nearly equal numbers. This is catastrophic cancellation, often driving the computed numerator to zero in FP32.

\subsection{Mixed and FP64 (stable form via Vieta)}
The mixed and FP64 programs compute the large root first:
\[
r_{\text{large}} = \frac{-b - \sqrt{\Delta}}{2a},
\]
and then use Vieta's product relation \(r_{\text{small}}r_{\text{large}} = c/a\):
\[
r_{\text{small}} = \frac{c}{a\,r_{\text{large}}}.
\]

This avoids the unstable subtraction \((-b + \sqrt{\Delta})\) and is numerically much more robust.

\section{Quantity 2: Variance of the Generated Data}

\subsection{FP32 baseline (one-pass formula)}
Define
\[
S_1 = \sum_{i=0}^{N-1} x_i,
\qquad
S_2 = \sum_{i=0}^{N-1} x_i^2,
\qquad
\mu = \frac{S_1}{N}.
\]
The FP32 baseline returns
\[
\sigma^2_{\text{one-pass}} = \frac{S_2}{N} - \mu^2.
\]

For large-magnitude data (here dominated by \(10^8\)), both terms are large and close, so subtracting them causes severe cancellation.

\subsection{Mixed and FP64 (two-pass formula)}
First pass:
\[
\mu = \frac{1}{N}\sum_{i=0}^{N-1} x_i.
\]
Second pass:
\[
\sigma^2_{\text{two-pass}} = \frac{1}{N}\sum_{i=0}^{N-1}(x_i-\mu)^2.
\]

This form is better conditioned for the centered residuals and avoids subtracting two large nearly equal moments.

\section{Algorithmic and Precision Changes from FP32 to Mixed}
The mixed version is not only a type cast; it combines selective precision promotion with algorithm replacement.

\begin{center}
\begin{tabular}{@{}lll@{}}
\toprule
Component & FP32 baseline & Mixed version \\
\midrule
Data storage \(x_i\) & FP32 (float) & FP32 (float, unchanged) \\
Quadratic root & Unstable \((-b+\sqrt{\Delta})/(2a)\) in FP32 & Stable Vieta form in FP64 \\
Variance & One-pass \(E[x^2]-E[x]^2\) in FP32 & Two-pass centered sum in FP64 \\
Accumulators & FP32 & FP64 \\
\bottomrule
\end{tabular}
\end{center}

\section{Why These Specific Changes}
FPChecker rounding-error reports in this example identify large relative errors at:
\begin{itemize}
  \item the quadratic cancellation numerator \((-b+\sqrt{\Delta})\), and
  \item the one-pass variance subtraction \((S_2/N - \mu^2)\).
\end{itemize}

Hence mixed precision targets exactly those hotspots:
\begin{enumerate}
  \item Promote sensitive arithmetic to FP64.
  \item Replace unstable formulas with numerically stable equivalents.
\end{enumerate}

Low-error operations (for example, data generation terms) remain FP32 to preserve the mixed-precision character.

\section{Expected Accuracy Behavior}
Relative to the FP64 reference:
\begin{itemize}
  \item FP32 often has very large relative error for \(r_{\text{small}}\) and for variance.
  \item Mixed should reduce these errors substantially.
\end{itemize}

In this example setup, the mixed result typically matches the FP64 quadratic small root closely and improves variance by many orders of magnitude compared with FP32.

\section{Build Instructions}
From the \texttt{tutorial/example\_6} directory:
\begin{verbatim}
pdflatex example6_math_formulas.tex
\end{verbatim}
If you prefer DVI first:
\begin{verbatim}
latex example6_math_formulas.tex
\end{verbatim}

\end{document}
